% TOP_PROBLEM: {"id": 30006689, "title": "Prove the Penrose inequality or present a counterexample", "category_id": 16, "category": {"id": 16, "name": "physics", "display_name": "Mathematical Physics", "description": "Problems at the intersection of mathematics and physics.", "slug": "physics", "order_index": 16, "created_at": "2026-07-31T15:26:25.671Z"}, "status": "open"}
% FIRST_POSTED: 2026-10-06
% ATTEMPT_STATUS: unresolved
% OPENAI_CLAIM: {"schema_version":1,"source_repo":"https://github.com/openai/math","source_commit":"adc7f1241b42e322a6451854ab7e4b4c146bf78a","release_date":"2026-10-06","match":"partial","resolution":"partial","independently_reviewed":false,"second_pass":"n/a","families":[{"family":"260","title":"Spacetime Penrose inequalities: enclosing area, charge, rotation, and anti-de Sitter extensions","subject":"Mathematical physics","manuscripts":[{"title":"Boundary graph deformations for the spacetime Penrose inequality","pdf_path":"preprints/Boundary-graph-deformations-for-the-spacetime-Penrose-inequality-September-27-2026/paper.pdf","readme_path":"preprints/Boundary-graph-deformations-for-the-spacetime-Penrose-inequality-September-27-2026/README.md","date":"2026-09-27","theorem_ref":"Theorem 2.4 (n3:intro:exterior-inequality) (neutral/00-setting.tex:177); Theorem 2.3 (n3:prepared:output) (neutral/00-setting.tex:124)","statement_sources":[{"path":"preprints/Boundary-graph-deformations-for-the-spacetime-Penrose-inequality-September-27-2026/build/neutral/00-setting.tex","line":177,"label":"Theorem 2.4 (n3:intro:exterior-inequality)"},{"path":"preprints/Boundary-graph-deformations-for-the-spacetime-Penrose-inequality-September-27-2026/build/neutral/00-setting.tex","line":124,"label":"Theorem 2.3 (n3:prepared:output)"}]},{"title":"Conformal flow and the Riemannian Penrose inequality with minimizing frontiers","pdf_path":"preprints/Conformal-flow-and-the-Riemannian-Penrose-inequality-with-minimizing-frontiers-September-27-2026/paper.pdf","readme_path":"preprints/Conformal-flow-and-the-Riemannian-Penrose-inequality-with-minimizing-frontiers-September-27-2026/README.md","date":"2026-09-27","theorem_ref":"Theorem 1.1 (thm:rpi) (sections/introduction.tex:120)","statement_sources":[{"path":"preprints/Conformal-flow-and-the-Riemannian-Penrose-inequality-with-minimizing-frontiers-September-27-2026/build/sections/introduction.tex","line":120,"label":"Theorem 1.1 (thm:rpi)"}]},{"title":"Equality and rigidity in the spacetime Penrose inequality","pdf_path":"preprints/Equality-and-rigidity-in-the-spacetime-Penrose-inequality-September-27-2026/paper.pdf","readme_path":"preprints/Equality-and-rigidity-in-the-spacetime-Penrose-inequality-September-27-2026/README.md","date":"2026-09-27","theorem_ref":"Theorem 2.5 (thm:rigidity) (sections/01-setting.tex:154); Theorem 2.6 (thm:converse) (sections/01-setting.tex:176); Theorem 8.2 (q3:intro:rigidity) (three/01-equality.tex:36); Theorem 8.3 (q3:intro:connected-rigidity) (three/01-equality.tex:79)","statement_sources":[{"path":"preprints/Equality-and-rigidity-in-the-spacetime-Penrose-inequality-September-27-2026/build/sections/01-setting.tex","line":154,"label":"Theorem 2.5 (thm:rigidity)"},{"path":"preprints/Equality-and-rigidity-in-the-spacetime-Penrose-inequality-September-27-2026/build/sections/01-setting.tex","line":176,"label":"Theorem 2.6 (thm:converse)"},{"path":"preprints/Equality-and-rigidity-in-the-spacetime-Penrose-inequality-September-27-2026/build/three/01-equality.tex","line":36,"label":"Theorem 8.2 (q3:intro:rigidity)"},{"path":"preprints/Equality-and-rigidity-in-the-spacetime-Penrose-inequality-September-27-2026/build/three/01-equality.tex","line":79,"label":"Theorem 8.3 (q3:intro:connected-rigidity)"}]},{"title":"Spacetime Penrose inequalities: enclosing area, charge, and rigidity","pdf_path":"preprints/Spacetime-Penrose-inequalities-enclosing-area-charge-and-rigidity-October-5-2026/paper.pdf","readme_path":"preprints/Spacetime-Penrose-inequalities-enclosing-area-charge-and-rigidity-October-5-2026/README.md","date":"2026-10-05","theorem_ref":"Theorem 2.2 (n3:intro:exterior-inequality) (sections/neutral/setup.tex:92); ch:intro:main (sections/charged/sections/02-preliminaries.tex:63); Corollary 9.4 (bridge:bridge:timelike) (sections/neutral-bridge.tex:99); Theorem 9.2 (bridge:intro:global-penrose) (sections/neutral-bridge.tex:42); Theorem 10.2 (n3:intro:rigidity) (sections/rigidity/sections/00-rigidity-setup.tex:63)","statement_sources":[{"path":"preprints/Spacetime-Penrose-inequalities-enclosing-area-charge-and-rigidity-October-5-2026/build/sections/neutral/setup.tex","line":92,"label":"Theorem 2.2 (n3:intro:exterior-inequality)"},{"path":"preprints/Spacetime-Penrose-inequalities-enclosing-area-charge-and-rigidity-October-5-2026/build/sections/charged/sections/02-preliminaries.tex","line":63,"label":"ch:intro:main"},{"path":"preprints/Spacetime-Penrose-inequalities-enclosing-area-charge-and-rigidity-October-5-2026/build/sections/neutral-bridge.tex","line":99,"label":"Corollary 9.4 (bridge:bridge:timelike)"},{"path":"preprints/Spacetime-Penrose-inequalities-enclosing-area-charge-and-rigidity-October-5-2026/build/sections/neutral-bridge.tex","line":42,"label":"Theorem 9.2 (bridge:intro:global-penrose)"},{"path":"preprints/Spacetime-Penrose-inequalities-enclosing-area-charge-and-rigidity-October-5-2026/build/sections/rigidity/sections/00-rigidity-setup.tex","line":63,"label":"Theorem 10.2 (n3:intro:rigidity)"}]},{"title":"The Kerr–Newman Penrose Inequality for Axisymmetric Electrovacuum Exteriors","pdf_path":"preprints/The-Kerr-Newman-Penrose-Inequality-for-Axisymmetric-Electrovacuum-Exteriors-October-5-2026/paper.pdf","readme_path":"preprints/The-Kerr-Newman-Penrose-Inequality-for-Axisymmetric-Electrovacuum-Exteriors-October-5-2026/README.md","date":"2026-10-05","theorem_ref":"Theorem 2.2 (m:main) (main.tex:293)","statement_sources":[{"path":"preprints/The-Kerr-Newman-Penrose-Inequality-for-Axisymmetric-Electrovacuum-Exteriors-October-5-2026/build/main.tex","line":293,"label":"Theorem 2.2 (m:main)"}]},{"title":"The spacetime Penrose inequality and enclosing area","pdf_path":"preprints/The-spacetime-Penrose-inequality-and-enclosing-area-September-27-2026/paper.pdf","readme_path":"preprints/The-spacetime-Penrose-inequality-and-enclosing-area-September-27-2026/README.md","date":"2026-09-27","theorem_ref":"Theorem 1.3 (thm:main) (sections/introduction.tex:87); Corollary 1.4 (cor:automatic-timelikeness) (sections/introduction.tex:110); Theorem 1.5 (thm:low) (sections/introduction.tex:120); Theorem 7.14 (bridge:bridge:timelike / thm:weak3) (sections/weak3.tex:1269); Theorem 7.16 (bridge:intro:global-penrose) (sections/weak3.tex:1354)","statement_sources":[{"path":"preprints/The-spacetime-Penrose-inequality-and-enclosing-area-September-27-2026/build/sections/introduction.tex","line":87,"label":"Theorem 1.3 (thm:main)"},{"path":"preprints/The-spacetime-Penrose-inequality-and-enclosing-area-September-27-2026/build/sections/introduction.tex","line":110,"label":"Corollary 1.4 (cor:automatic-timelikeness)"},{"path":"preprints/The-spacetime-Penrose-inequality-and-enclosing-area-September-27-2026/build/sections/introduction.tex","line":120,"label":"Theorem 1.5 (thm:low)"},{"path":"preprints/The-spacetime-Penrose-inequality-and-enclosing-area-September-27-2026/build/sections/weak3.tex","line":1269,"label":"Theorem 7.14 (bridge:bridge:timelike / thm:weak3)"},{"path":"preprints/The-spacetime-Penrose-inequality-and-enclosing-area-September-27-2026/build/sections/weak3.tex","line":1354,"label":"Theorem 7.16 (bridge:intro:global-penrose)"}]},{"title":"The spacetime Penrose inequality with charge and original-data rigidity","pdf_path":"preprints/The-spacetime-Penrose-inequality-with-charge-and-original-data-rigidity-October-5-2026/paper.pdf","readme_path":"preprints/The-spacetime-Penrose-inequality-with-charge-and-original-data-rigidity-October-5-2026/README.md","date":"2026-10-05","theorem_ref":"Theorem 2.3 (thm:numerical) (sections/statements.tex:103); Theorem 2.5 (thm:rigidity) (sections/statements.tex:139); Theorem 2.6 (thm:converse) (sections/statements.tex:183)","statement_sources":[{"path":"preprints/The-spacetime-Penrose-inequality-with-charge-and-original-data-rigidity-October-5-2026/build/sections/statements.tex","line":103,"label":"Theorem 2.3 (thm:numerical)"},{"path":"preprints/The-spacetime-Penrose-inequality-with-charge-and-original-data-rigidity-October-5-2026/build/sections/statements.tex","line":139,"label":"Theorem 2.5 (thm:rigidity)"},{"path":"preprints/The-spacetime-Penrose-inequality-with-charge-and-original-data-rigidity-October-5-2026/build/sections/statements.tex","line":183,"label":"Theorem 2.6 (thm:converse)"}]}],"lean":{"doc_path":"lean/docs/260.md","comparators":[],"declarations":[],"covers_main_theorem":false},"reasoning_trace":null}]}
\documentclass[11pt]{article}
\usepackage[margin=25mm]{geometry}
\usepackage{fontspec}
\setmainfont{Latin Modern Roman}
\usepackage{amsmath,amssymb,mathtools}
\usepackage{unicode-math}
\setmathfont{Latin Modern Math}
\usepackage{enumitem}
\usepackage{xurl,hyperref,bookmark}
\hypersetup{colorlinks=true,urlcolor=blue}
\setcounter{secnumdepth}{0}
\setlength{\parindent}{0pt}
\setlength{\parskip}{5pt}
\setlength{\emergencystretch}{3em}
\begin{document}
\section*{Prove the Penrose inequality or present a counterexample}

\subsection{OpenAI's claim}
The linked manuscripts present the following claims by OpenAI: In dimensions three and four, trapped exteriors are claimed deformable to nonnegative-scalar-curvature metrics with protected enclosing area and arbitrarily small ADM-energy loss; The numerical Riemannian mass-area inequality is claimed in every dimension for full outer-minimizing, locally perimeter-minimizing frontiers, including singular and disconnected ones; Equality in the neutral spacetime bound is claimed to reconstruct the original exterior as a Schwarzschild--Tangherlini slice under the precise horizon and decay assumptions; A combined manuscript claims neutral three- and four-dimensional bounds, three-dimensional charged upper-area bounds, original-data rigidity, and a local anti-de Sitter perturbation inequality; For the designated axisymmetric electrovacuum class, the claimed bound is $m^2\geq A/(16\pi)+Q^2/2+\pi(Q^4+4J^2)/A$ with conserved total angular momentum; The proposed neutral spacetime Penrose inequality compares invariant ADM mass with minimum enclosing area in all dimensions $n\geq3$, under the listed energy, decay, and weak future-trapping hypotheses; For one-ended three-dimensional source-free electric--magnetic data, the proposed sharp bound is $m\geq Q$ and $r_A\leq m+\sqrt{m^2-Q^2}$, with Reissner--Nordstrom rigidity in the stated horizon subclass. Together they provide the following support within the scopes described below.

\subsection{Mathematical statement of the claimed result}
\paragraph{Result 1.} Let $\Omega^n$, $n\ge3$, be smooth, connected and orientable, with nonempty compact smooth boundary $S$, complete with that boundary included, and exactly one end, with Euclidean coordinates and compact complement. The tensors $g,K$ extend smoothly to $S$. With $\tau=\operatorname{tr}_gK$, use
\[
 2\mu=R_g+\tau^2-|K|_g^2,\qquad J=\operatorname{div}_g(K-\tau g),\qquad
 \mu\ge|J|_g,\quad\mu,|J|_g\in L^1(\Omega,dV_g).
\]
The finite ADM fluxes use $\omega_{n-1}=|\mathbb S^{n-1}|$ and
\[
 E=\frac1{2(n-1)\omega_{n-1}}\lim_{R\to\infty}\int_{|x|=R}
 (\partial_jg_{ij}-\partial_ig_{jj})n_\delta^i\,dA_\delta,\quad
 P_i=\frac1{(n-1)\omega_{n-1}}\lim_{R\to\infty}\int_{|x|=R}(K_{ij}-\tau g_{ij})n_\delta^j\,dA_\delta.
\]
Orient boundary normals towards the end, with $\theta_+=H+\operatorname{tr}_SK$. Here $O_j(r^{-b})$ controls coordinate derivatives through order $j$, of respective sizes $O(r^{-b-i})$. Choose a connected smooth end-side region $D\subset\Omega$, closed in $\Omega$, containing the distant end and having its manifold interior in $\operatorname{int}\Omega$. A full cut consists of all of its intrinsic boundary $\Gamma=\partial D$; require this boundary to be compact, smoothly embedded and two-sided. All components and portions coinciding with $S=\partial\Omega$ count. Put $a_g(S)=\inf_\Gamma\operatorname{Area}_g(\Gamma)$. In dimension three, suppose $g-\delta=O_2(r^{-q})$, $K=O_1(r^{-1-q})$ for $q>1/2$, and $\theta_+\le0$ on every boundary component. If $E>|P|$, the exterior theorem asserts
\[
 \sqrt{E^2-|P|^2}\ge\sqrt{a_g(S)/(16\pi)}.
\]
The preparatory output treats the stronger subclass with compactly supported $K$, $g-\delta=O_j(r^{-1})$ for every fixed $j$, $P=0$, $E>0$, $\theta_+<0$, $A_*=a_g(S)>0$, and, for a smooth radius extension $\varrho\ge1$,
\[
 R_g=O(r^{-3-\delta}),\qquad \mu-|J|_g\ge c\varrho^{-3-\delta}\quad(c>0,\ 0<\delta<1).
\]
For sufficiently small fixed $\epsilon>0$ and large $N$, a smooth one-ended subexterior has a complete boundary-included metric $\widehat g_{\epsilon,N}$ satisfying
\[
 R_{\widehat g}\ge0,\quad H_{\partial\Omega}(\widehat g)<0,\quad
 \widehat g\ge e^{-4\epsilon}g,\quad E_{\widehat g}\le E+\eta_{\epsilon,N},\quad
 |T|_{\widehat g}\ge e^{-4\epsilon}A_*
\]
for every full cut $T$. At fixed $\epsilon$, $0\le\eta_{\epsilon,N}\le C(1+N)^a(e^{-\epsilon N}+e^{-\epsilon N/2})\to0$. The output metric has $O_2(r^{-1})$ decay and scalar curvature $O(r^{-3-\delta'})$ for some $\delta'>0$.

\paragraph{Result 2.} For $n\ge3$, let $(E,h)$ be a connected complete exterior with one Euclidean end and compact complement, including its complete possibly singular frontier $\Sigma$. Completeness is measured within each finite-distance class of its intrinsic length metric. The metric extends smoothly through the frontier and satisfies
\[
 R_h\ge0,\quad R_h\in L^1(E,dV_h),\quad R_h=O(r^{-n-\beta})\ (\beta>0),\qquad
 h-\delta=O_2(r^{-\tau})\ (\tau>(n-2)/2).
\]
The full frontier is the perimeter support of a filled side, outer minimizing and locally perimeter minimizing on both sides. Its regular part is smooth, embedded and minimal; its singular part satisfies $\dim_{\mathrm H}\Sigma_{\mathrm{sing}}\le n-8$. With $\omega_{n-1}=|\mathbb S^{n-1}|$ and the usual ADM energy normalization $[2(n-1)\omega_{n-1}]^{-1}$,
\[
 m_{\mathrm{ADM}}(h)\ge\frac12\left(\frac{\mathcal H_h^{n-1}(\Sigma)}{\omega_{n-1}}\right)^{(n-2)/(n-1)}.
\]
Every frontier component contributes to the area; no spin or frontier-connectedness hypothesis is imposed.

\paragraph{Result 3.} Let $\Omega^n$, $n\ge3$, be smooth, connected and orientable, with nonempty compact smooth boundary $S$, complete with that boundary included, and exactly one end, with Euclidean coordinates and compact complement. The tensors $g,K$ extend smoothly to $S$. With $\tau=\operatorname{tr}_gK$, use
\[
 2\mu=R_g+\tau^2-|K|_g^2,\qquad J=\operatorname{div}_g(K-\tau g),\qquad
 \mu\ge|J|_g,\quad\mu,|J|_g\in L^1(\Omega,dV_g).
\]
The finite ADM fluxes use $\omega_{n-1}=|\mathbb S^{n-1}|$ and
\[
 E=\frac1{2(n-1)\omega_{n-1}}\lim_{R\to\infty}\int_{|x|=R}
 (\partial_jg_{ij}-\partial_ig_{jj})n_\delta^i\,dA_\delta,\quad
 P_i=\frac1{(n-1)\omega_{n-1}}\lim_{R\to\infty}\int_{|x|=R}(K_{ij}-\tau g_{ij})n_\delta^j\,dA_\delta.
\]
Orient boundary normals towards the end, with $\theta_+=H+\operatorname{tr}_SK$. Here $O_j(r^{-b})$ controls coordinate derivatives through order $j$, of respective sizes $O(r^{-b-i})$. Choose a connected smooth end-side region $D\subset\Omega$, closed in $\Omega$, containing the distant end and having its manifold interior in $\operatorname{int}\Omega$. A full cut consists of all of its intrinsic boundary $\Gamma=\partial D$; require this boundary to be compact, smoothly embedded and two-sided. All components and portions coinciding with $S=\partial\Omega$ count. Put $a_g(S)=\inf_\Gamma\operatorname{Area}_g(\Gamma)$. Suppose $g-\delta=O_6(r^{-q})$, $K=O_5(r^{-1-q})$, $(n-2)/2<q<n-2$, and $E>|P|$. Require connected $S$, $\theta_+(S)=0$, and $\operatorname{Area}_g(\Gamma)\ge A=|S|_g>0$ for every full cut. Exclude every compact smooth embedded two-sided full enclosing hypersurface wholly in $\operatorname{int}\Omega$ with $\theta_+\le0$, allowing disconnected excluded surfaces. If
\[
 m=\sqrt{E^2-|P|^2}=\frac12(A/\omega_{n-1})^{(n-2)/(n-1)},
\]
then the entire original pair $(g,K)$, including the boundary and future normal, is induced by a smooth spacelike embedding into Schwarzschild--Tangherlini of mass $m$, using its regular extension across the future horizon. The interior lies in its domain of outer communication, the end reaches spatial infinity, and $S$ maps to a full future-horizon section or the bifurcation sphere. Conversely, admissible spacelike model slices with these same intrinsic hypotheses attain equality and have $A=\omega_{n-1}(2m)^{(n-1)/(n-2)}$. Time-symmetric and nonzero-$K$ examples exist in each $n\ge3$. In the three-dimensional complete-data equality theorem, use a smooth connected orientable boundaryless $M$ with complete metric, compact core and finitely many AF ends, each obeying $O_2/O_1$ decay with $q>1/2$, finite ADM limits, integrable densities and DEC. Take $S=\bigsqcup_{i=1}^{\ell}S_i$, $1\le\ell<\infty$, with disjoint connected compact smooth two-sided components bounding a connected closed exterior $\Omega$ with exactly the designated end. Require $\theta_+(S_i)=0$ and $\operatorname{Area}_g(\Gamma)\ge A=\sum_i|S_i|_g>0$ for every full enclosing cut. Exclude distinct enclosing cuts whose components each coincide with an entire $S_i$ or lie wholly outside, with at least one component of the latter kind, and with $\theta_+\le0$ on all components. If $E>|P|$ and $\sqrt{E^2-|P|^2}=\sqrt{A/(16\pi)}$, then $\ell=1$, $S\cong S^2$, and the entire original exterior admits a proper smooth spacelike Schwarzschild embedding of mass $\sqrt{E^2-|P|^2}$ inducing $g,K$ and the prescribed future normal. Its end reaches spatial infinity and its boundary is a full future-horizon section or bifurcation sphere, with smooth continuation through that boundary. For connected $S$, the connected-horizon corollary replaces the cut exclusion by exclusion of every wholly exterior weakly future-trapped enclosing surface, including disconnected ones. Thus $\ell\ge2$ forces strict inequality in the finite-component subclass.

\paragraph{Result 4.} Let $\Omega^n$, $n\ge3$, be smooth, connected and orientable, with nonempty compact smooth boundary $S$, complete with that boundary included, and exactly one end, with Euclidean coordinates and compact complement. The tensors $g,K$ extend smoothly to $S$. With $\tau=\operatorname{tr}_gK$, use
\[
 2\mu=R_g+\tau^2-|K|_g^2,\qquad J=\operatorname{div}_g(K-\tau g),\qquad
 \mu\ge|J|_g,\quad\mu,|J|_g\in L^1(\Omega,dV_g).
\]
The finite ADM fluxes use $\omega_{n-1}=|\mathbb S^{n-1}|$ and
\[
 E=\frac1{2(n-1)\omega_{n-1}}\lim_{R\to\infty}\int_{|x|=R}
 (\partial_jg_{ij}-\partial_ig_{jj})n_\delta^i\,dA_\delta,\quad
 P_i=\frac1{(n-1)\omega_{n-1}}\lim_{R\to\infty}\int_{|x|=R}(K_{ij}-\tau g_{ij})n_\delta^j\,dA_\delta.
\]
Orient boundary normals towards the end, with $\theta_+=H+\operatorname{tr}_SK$. Here $O_j(r^{-b})$ controls coordinate derivatives through order $j$, of respective sizes $O(r^{-b-i})$. Choose a connected smooth end-side region $D\subset\Omega$, closed in $\Omega$, containing the distant end and having its manifold interior in $\operatorname{int}\Omega$. A full cut consists of all of its intrinsic boundary $\Gamma=\partial D$; require this boundary to be compact, smoothly embedded and two-sided. All components and portions coinciding with $S=\partial\Omega$ count. Put $a_g(S)=\inf_\Gamma\operatorname{Area}_g(\Gamma)$. In dimension three, suppose $g-\delta=O_2(r^{-q})$, $K=O_1(r^{-1-q})$ for $q>1/2$, and $\theta_+\le0$ on every boundary component. The neutral result gives $\sqrt{E^2-|P|^2}\ge\sqrt{a_g(S)/(16\pi)}$ when $E>|P|$, and a separate endpoint argument proves $E>|P|$ without imposing that sign initially. In dimension three on the same smooth one-ended complete exterior, let the electric and magnetic fields $\mathcal E,\mathcal B$ be smooth, divergence free, and $O_1(r^{-2})$. Suppose $g-\delta=O_2(r^{-q})$, $K=O_1(r^{-1-q})$ for $q>1/2$, finite ADM and charge fluxes, and integrable geometric densities $2\mu=R_g+\tau^2-|K|^2$, $J=\operatorname{div}_g(K-\tau g)$. The non-electromagnetic matter condition is
\[
 \mu_m=\mu-|\mathcal E|^2-|\mathcal B|^2\ge|J_m|,\qquad
 J_m=J-2(\mathcal E\times\mathcal B)^\flat.
\]
Set $Q_E=(4\pi)^{-1}\lim_R\int_{|x|=R}g(\mathcal E,\nu_R)dA_g$, define $Q_B$ similarly, and put $Q=\sqrt{Q_E^2+Q_B^2}$, $r_A=\sqrt{a_g(S)/(4\pi)}$. If $\theta_+\le0$ on the entire boundary, the charged conclusion is
\[
 E>|P|,\qquad m=\sqrt{E^2-|P|^2}\ge Q,\qquad
 r_A\le m+\sqrt{m^2-Q^2}.
\]
For $r_A>Q$ this is equivalent to $m\ge(r_A+Q^2/r_A)/2$; for $0<r_A\le Q$, the mass-charge bound is the active restriction. The combined paper's numerical charged theorem does not assert unrestricted equality rigidity. There is also a complete-data formulation on a smooth connected orientable boundaryless three-manifold $M$, with complete $g$, a compact core and finitely many AF ends. Each end has $g-\delta=O_2(r^{-q})$, $K=O_1(r^{-1-q})$, $q>1/2$, and finite ADM limits; the constraint densities are globally integrable and satisfy DEC. Choose an end-side open region $D^+$, possibly disconnected or containing other ends, whose full boundary $S$ is nonempty, compact, smooth, embedded and two-sided, with $\theta_+\le0$ towards $D^+$. Define
\[
 A_{\min}(S)=\inf_{D'^+\subseteq D^+}\operatorname{Area}_g(\partial D'^+),
\]
where smooth open competitors contain the distant designated end and have compact smooth full boundary, permitting disconnected competitors, other retained ends, and coincidence with $S$. The conclusion is $E>|P|$ and $\sqrt{E^2-|P|^2}\ge\sqrt{A_{\min}(S)/(16\pi)}$. In the three-dimensional complete-data equality theorem, use a smooth connected orientable boundaryless $M$ with complete metric, compact core and finitely many AF ends, each obeying $O_2/O_1$ decay with $q>1/2$, finite ADM limits, integrable densities and DEC. Take $S=\bigsqcup_{i=1}^{\ell}S_i$, $1\le\ell<\infty$, with disjoint connected compact smooth two-sided components bounding a connected closed exterior $\Omega$ with exactly the designated end. Require $\theta_+(S_i)=0$ and $\operatorname{Area}_g(\Gamma)\ge A=\sum_i|S_i|_g>0$ for every full enclosing cut. Exclude distinct enclosing cuts whose components each coincide with an entire $S_i$ or lie wholly outside, with at least one component of the latter kind, and with $\theta_+\le0$ on all components. If $E>|P|$ and $\sqrt{E^2-|P|^2}=\sqrt{A/(16\pi)}$, then $\ell=1$, $S\cong S^2$, and the entire original exterior admits a proper smooth spacelike Schwarzschild embedding of mass $\sqrt{E^2-|P|^2}$ inducing $g,K$ and the prescribed future normal. Its end reaches spatial infinity and its boundary is a full future-horizon section or bifurcation sphere, with smooth continuation through that boundary. For connected $S$, the connected-horizon corollary replaces the cut exclusion by exclusion of every wholly exterior weakly future-trapped enclosing surface, including disconnected ones. Thus $\ell\ge2$ forces strict inequality in the finite-component subclass.

\paragraph{Result 5.} Let the smooth oriented connected exterior be diffeomorphic to \(\mathbb R^3\) minus an open ball, with spherical smooth compact boundary \(S\), complete metric including \(S\), a single AF end, and smooth \(g,K,\mathcal E,\mathcal B\). A smooth effective \(U(1)\) action preserves all data and the boundary, with rotation generator of period \(2\pi\) agreeing with ordinary axial rotation at infinity. Impose source-free electrovacuum constraints
\[
16\pi\mu=R_g+(\operatorname{tr}_gK)^2-|K|^2,
\quad 8\pi J_{\rm con}=\operatorname{div}_g(K-(\operatorname{tr}_gK)g),
\]
\[
8\pi\mu=|\mathcal E|^2+|\mathcal B|^2,\quad
4\pi J_{\rm con}=(\mathcal E\times\mathcal B)^\flat,\quad
\operatorname{div}_g\mathcal E=\operatorname{div}_g\mathcal B=0.
\]
Use \(g-\delta=O_4(r^{-1})\), \(K=O_3(r^{-3})\), Coulomb leading fields
\(\mathcal E^i=Q_e x^i/r^3+O_3(r^{-3})\), \(\mathcal B^i=Q_b x^i/r^3+O_3(r^{-3})\), integrable constraints and finite ADM energy, momentum, angular momentum and charge fluxes. Require \(P_{\rm ADM}=0\), \(m=E_{\rm ADM}>0\), \(Q^2=Q_e^2+Q_b^2\), and let \(J\) denote conserved angular momentum, with its electromagnetic boundary contribution included; it agrees with \(J_{\rm ADM}\). The boundary is a future MOTS; no wholly exterior enclosing smooth surface is weakly outer trapped; \(\operatorname{Area}_g(\Gamma)\ge A=|S|_g>0\) for every full cut \(\Gamma\); and
\[A\ge4\pi\sqrt{Q^4+4J^2}.\]
\[
m^2\ge\frac{A}{16\pi}+\frac{Q^2}{2}+\frac{\pi(Q^4+4J^2)}{A}.
\]
Equality on the strict area branch \(A>4\pi\sqrt{Q^4+4J^2}\) identifies the entire original data with a smooth spacelike slice of subextremal dyonic Kerr--Newman with matching mass, angular momentum and charges, including smooth attachment and future normal at the horizon and an end at spatial infinity. Conversely, admissible such slices attain equality. The numerical endpoint is included, but extremal equality is not classified. Here every full cut is the entire boundary of an end-side exterior region, including all boundary-coincident components, and $O_j$ includes coordinate derivatives through order $j$.

\paragraph{Result 6.} Let $\Omega^n$, $n\ge3$, be smooth, connected and orientable, with nonempty compact smooth boundary $S$, complete with that boundary included, and exactly one end, with Euclidean coordinates and compact complement. The tensors $g,K$ extend smoothly to $S$. With $\tau=\operatorname{tr}_gK$, use
\[
 2\mu=R_g+\tau^2-|K|_g^2,\qquad J=\operatorname{div}_g(K-\tau g),\qquad
 \mu\ge|J|_g,\quad\mu,|J|_g\in L^1(\Omega,dV_g).
\]
The finite ADM fluxes use $\omega_{n-1}=|\mathbb S^{n-1}|$ and
\[
 E=\frac1{2(n-1)\omega_{n-1}}\lim_{R\to\infty}\int_{|x|=R}
 (\partial_jg_{ij}-\partial_ig_{jj})n_\delta^i\,dA_\delta,\quad
 P_i=\frac1{(n-1)\omega_{n-1}}\lim_{R\to\infty}\int_{|x|=R}(K_{ij}-\tau g_{ij})n_\delta^j\,dA_\delta.
\]
Orient boundary normals towards the end, with $\theta_+=H+\operatorname{tr}_SK$. Here $O_j(r^{-b})$ controls coordinate derivatives through order $j$, of respective sizes $O(r^{-b-i})$. Choose a connected smooth end-side region $D\subset\Omega$, closed in $\Omega$, containing the distant end and having its manifold interior in $\operatorname{int}\Omega$. A full cut consists of all of its intrinsic boundary $\Gamma=\partial D$; require this boundary to be compact, smoothly embedded and two-sided. All components and portions coinciding with $S=\partial\Omega$ count. Put $a_g(S)=\inf_\Gamma\operatorname{Area}_g(\Gamma)$. For $g-\delta=O_6(r^{-q})$, $K=O_5(r^{-1-q})$ with $(n-2)/2<q<n-2$, $\theta_+\le0$, $a_g(S)>0$ and $E>|P|$, the numerical bound is
\[
 \sqrt{E^2-|P|^2}\ge\frac12\left(\frac{a_g(S)}{\omega_{n-1}}\right)^{(n-2)/(n-1)}.
\]
Its endpoint corollary derives $E>|P|$ from the other hypotheses. In $n\in\{3,4\}$, the lower-regularity theorem only requires $g-\delta=O_2(r^{-q})$, $K=O_1(r^{-1-q})$ with $q>(n-2)/2$; it proves $a_g(S)>0$ and $E>|P|$ as well as the same inequality. The coefficient is sharp. Numerical assumptions do not require connected boundary, maximality, vacuum, outermostness, spin, or outer area minimization. There is also a complete-data formulation on a smooth connected orientable boundaryless three-manifold $M$, with complete $g$, a compact core and finitely many AF ends. Each end has $g-\delta=O_2(r^{-q})$, $K=O_1(r^{-1-q})$, $q>1/2$, and finite ADM limits; the constraint densities are globally integrable and satisfy DEC. Choose an end-side open region $D^+$, possibly disconnected or containing other ends, whose full boundary $S$ is nonempty, compact, smooth, embedded and two-sided, with $\theta_+\le0$ towards $D^+$. Define
\[
 A_{\min}(S)=\inf_{D'^+\subseteq D^+}\operatorname{Area}_g(\partial D'^+),
\]
where smooth open competitors contain the distant designated end and have compact smooth full boundary, permitting disconnected competitors, other retained ends, and coincidence with $S$. The conclusion is $E>|P|$ and $\sqrt{E^2-|P|^2}\ge\sqrt{A_{\min}(S)/(16\pi)}$.

\paragraph{Result 7.} Let $\Omega^n$, $n\ge3$, be smooth, connected and orientable, with nonempty compact smooth boundary $S$, complete with that boundary included, and exactly one end, with Euclidean coordinates and compact complement. The tensors $g,K$ extend smoothly to $S$. With $\tau=\operatorname{tr}_gK$, use
\[
 2\mu=R_g+\tau^2-|K|_g^2,\qquad J=\operatorname{div}_g(K-\tau g),\qquad
 \mu\ge|J|_g,\quad\mu,|J|_g\in L^1(\Omega,dV_g).
\]
The finite ADM fluxes use $\omega_{n-1}=|\mathbb S^{n-1}|$ and
\[
 E=\frac1{2(n-1)\omega_{n-1}}\lim_{R\to\infty}\int_{|x|=R}
 (\partial_jg_{ij}-\partial_ig_{jj})n_\delta^i\,dA_\delta,\quad
 P_i=\frac1{(n-1)\omega_{n-1}}\lim_{R\to\infty}\int_{|x|=R}(K_{ij}-\tau g_{ij})n_\delta^j\,dA_\delta.
\]
Orient boundary normals towards the end, with $\theta_+=H+\operatorname{tr}_SK$. Here $O_j(r^{-b})$ controls coordinate derivatives through order $j$, of respective sizes $O(r^{-b-i})$. Choose a connected smooth end-side region $D\subset\Omega$, closed in $\Omega$, containing the distant end and having its manifold interior in $\operatorname{int}\Omega$. A full cut consists of all of its intrinsic boundary $\Gamma=\partial D$; require this boundary to be compact, smoothly embedded and two-sided. All components and portions coinciding with $S=\partial\Omega$ count. Put $a_g(S)=\inf_\Gamma\operatorname{Area}_g(\Gamma)$. In dimension three on the same smooth one-ended complete exterior, let the electric and magnetic fields $\mathcal E,\mathcal B$ be smooth, divergence free, and $O_1(r^{-2})$. Suppose $g-\delta=O_2(r^{-q})$, $K=O_1(r^{-1-q})$ for $q>1/2$, finite ADM and charge fluxes, and integrable geometric densities $2\mu=R_g+\tau^2-|K|^2$, $J=\operatorname{div}_g(K-\tau g)$. The non-electromagnetic matter condition is
\[
 \mu_m=\mu-|\mathcal E|^2-|\mathcal B|^2\ge|J_m|,\qquad
 J_m=J-2(\mathcal E\times\mathcal B)^\flat.
\]
Set $Q_E=(4\pi)^{-1}\lim_R\int_{|x|=R}g(\mathcal E,\nu_R)dA_g$, define $Q_B$ similarly, and put $Q=\sqrt{Q_E^2+Q_B^2}$, $r_A=\sqrt{a_g(S)/(4\pi)}$. Each boundary component may satisfy either $H+\operatorname{tr}_SK\le0$ or $H-\operatorname{tr}_SK\le0$, with the choice made separately. The numerical assertion is
\[
 a_g(S)>0,\quad E>|P|,\quad m=\sqrt{E^2-|P|^2}>0,\quad
 E\ge\sqrt{|P|^2+Q^2},\quad m\ge Q,\quad r_A\le m+\sqrt{m^2-Q^2}.
\]
The polynomial expression $m\ge(r_A+Q^2/r_A)/2$ applies on $r_A>Q$; at $0<r_A\le Q$ the active condition is $m\ge Q$. In the connected future-horizon subclass, impose $\theta_+=0$, outer area minimization and exclusion of every wholly exterior weakly trapped full cut. If $m>Q$ and $\sqrt{A/(4\pi)}=m+\sqrt{m^2-Q^2}$, the entire original data, including fields, arise from a smooth spacelike slice of dyonic Reissner--Nordstrom with parameters $(m,Q_E,Q_B)$ and smooth future-horizon attachment. Matter densities vanish on that exterior. Conversely admissible model slices with actual mass $M>\sqrt{q_E^2+q_B^2}$ and charges $(q_E,q_B)$ attain
\[
 a_g(S)=|S|_g=4\pi\left(M+\sqrt{M^2-q_E^2-q_B^2}\right)^2.
\]
Extremal equality is not classified.

\subsection{Outline of the argument}
Collars and weak boundary supports set the graph barrier. Elliptic estimates exclude contact with the floor and solve the graph-conformal deformation; its negative inner mean curvature supplies the boundary route to the neutral mass-area bound. A conformal flow is built from smooth obstacles while preserving frontier area. Regular-patch estimates control forcing, and distance-median separation handles the singular region; mass monotonicity and the limiting comparison give the stated numerical inequality. Minimizing enclosures are treated through their area envelope. Saturation yields a constrained minimum in the original data, whose static rigidity equations produce the spacetime embedding and smooth horizon attachment; in three dimensions the analysis also rules out equality for disconnected horizons.

Prepared-end graph deformations protect the minimum area and compare invariant masses. Conformal flow supplies the Riemannian step, electromagnetic flux converts charge into the scalar comparison, and equality conditions reconstruct the original static spacetime data; the anti-de Sitter result uses a separate local conformal expansion. Axial variables are compared with a Kerr--Newman family at prescribed conformal capacity. Riemannian scale comparison and electromagnetic angular-momentum bookkeeping give the mass bound; saturation on the strict physical-area branch reconstructs an admissible dyonic spacetime slice. End preparation removes ADM momentum before a coupled graph and conformal scalar construction. Filled-equation estimates and height separation protect enclosing cuts, while the scalar identity bounds the change in mass; the Riemannian comparison then returns the inequality to the original data.

A coupled filled scalar deformation retains enclosing area while accounting for electromagnetic flux in the mass balance. The neutral comparison gives the upper-area branch, and saturation of the comparison identities reconstructs a dyonic spacetime slice carrying the given metric and extrinsic curvature.

\subsection{Scope relative to this problem}
The cited horizon formulation compares ADM energy $E_{\mathrm{ADM}}$ with the actual horizon area $A_\Sigma$. The selected OpenAI papers claim the distinct minimum-enclosing-area inequality $\sqrt{E^2-|P|^2}\ge\sqrt{A_{\min}/(16\pi)}$, including weak-decay Theorem 1.5 with $g-\delta=O_2(r^{-q})$, $K=O_1(r^{-1-q})$, $q>1/2$, in dimension three. That theorem also derives $A_{\min}>0$ and $E>|P|$. Although $E\ge\sqrt{E^2-|P|^2}$, the same implication does not replace $A_{\min}$ by the potentially larger $A_\Sigma$. The original-data Schwarzschild equality claims impose outer area minimization and additional connected-horizon or finite-component outermostness conditions. These results give a positive partial solution through the enclosing-area numerical theorem and the specified rigidity subclasses; they do not cover the unrestricted raw-horizon-area and equality formulation.

\subsection{Formal verification}
For family 260, the only selected family entry is the CKS end-replacement construction and Schwarzschild equality examples. A connected-boundary inequality is conditional on the flat Penrose input; the general Bondi or spacetime inequalities are not formalized there. No Comparator is assigned to the selected manuscript here; the family page selects other papers or supporting statements. This site has not executed the Lean code. The scope descriptions do not certify the remaining manuscript arguments.

\subsection{Status caveat}
OpenAI attributes these manuscripts to work by an internal model. Its release cautions that arguments without formalization can contain mistakes. This catalogue has not independently checked the proofs; the labels report the mathematical claims attributed to OpenAI.

\subsection{Sources}
\begin{itemize}
\item OpenAI, \emph{Boundary graph deformations for the spacetime Penrose inequality}, 2026-09-27; citation key \texttt{OAI:Boundary-graph-deformations-for-the-spacetime-Penrose-inequality-September-27-2026}. \href{https://github.com/openai/math/blob/adc7f1241b42e322a6451854ab7e4b4c146bf78a/preprints/Boundary-graph-deformations-for-the-spacetime-Penrose-inequality-September-27-2026/paper.pdf}{PDF (GitHub)}; \href{https://github.com/openai/math/raw/adc7f1241b42e322a6451854ab7e4b4c146bf78a/preprints/Boundary-graph-deformations-for-the-spacetime-Penrose-inequality-September-27-2026/paper.pdf}{PDF (direct)}; \href{https://github.com/openai/math/blob/adc7f1241b42e322a6451854ab7e4b4c146bf78a/preprints/Boundary-graph-deformations-for-the-spacetime-Penrose-inequality-September-27-2026/README.md}{Manuscript page and citation}; \href{https://github.com/openai/math/blob/main/preprints/Boundary-graph-deformations-for-the-spacetime-Penrose-inequality-September-27-2026/paper.pdf}{Latest version}; Statement source: \href{https://github.com/openai/math/blob/adc7f1241b42e322a6451854ab7e4b4c146bf78a/preprints/Boundary-graph-deformations-for-the-spacetime-Penrose-inequality-September-27-2026/build/neutral/00-setting.tex\#L177}{Theorem 2.4 (n3:intro:exterior-inequality)}; Statement source: \href{https://github.com/openai/math/blob/adc7f1241b42e322a6451854ab7e4b4c146bf78a/preprints/Boundary-graph-deformations-for-the-spacetime-Penrose-inequality-September-27-2026/build/neutral/00-setting.tex\#L124}{Theorem 2.3 (n3:prepared:output)}.
\item OpenAI, \emph{Conformal flow and the Riemannian Penrose inequality with minimizing frontiers}, 2026-09-27; citation key \texttt{OAI:Conformal-flow-and-the-Riemannian-Penrose-inequality-with-minimizing-frontiers-September-27-2026}. \href{https://github.com/openai/math/blob/adc7f1241b42e322a6451854ab7e4b4c146bf78a/preprints/Conformal-flow-and-the-Riemannian-Penrose-inequality-with-minimizing-frontiers-September-27-2026/paper.pdf}{PDF (GitHub)}; \href{https://github.com/openai/math/raw/adc7f1241b42e322a6451854ab7e4b4c146bf78a/preprints/Conformal-flow-and-the-Riemannian-Penrose-inequality-with-minimizing-frontiers-September-27-2026/paper.pdf}{PDF (direct)}; \href{https://github.com/openai/math/blob/adc7f1241b42e322a6451854ab7e4b4c146bf78a/preprints/Conformal-flow-and-the-Riemannian-Penrose-inequality-with-minimizing-frontiers-September-27-2026/README.md}{Manuscript page and citation}; \href{https://github.com/openai/math/blob/main/preprints/Conformal-flow-and-the-Riemannian-Penrose-inequality-with-minimizing-frontiers-September-27-2026/paper.pdf}{Latest version}; Statement source: \href{https://github.com/openai/math/blob/adc7f1241b42e322a6451854ab7e4b4c146bf78a/preprints/Conformal-flow-and-the-Riemannian-Penrose-inequality-with-minimizing-frontiers-September-27-2026/build/sections/introduction.tex\#L120}{Theorem 1.1 (thm:rpi)}.
\item OpenAI, \emph{Equality and rigidity in the spacetime Penrose inequality}, 2026-09-27; citation key \texttt{OAI:Equality-and-rigidity-in-the-spacetime-Penrose-inequality-September-27-2026}. \href{https://github.com/openai/math/blob/adc7f1241b42e322a6451854ab7e4b4c146bf78a/preprints/Equality-and-rigidity-in-the-spacetime-Penrose-inequality-September-27-2026/paper.pdf}{PDF (GitHub)}; \href{https://github.com/openai/math/raw/adc7f1241b42e322a6451854ab7e4b4c146bf78a/preprints/Equality-and-rigidity-in-the-spacetime-Penrose-inequality-September-27-2026/paper.pdf}{PDF (direct)}; \href{https://github.com/openai/math/blob/adc7f1241b42e322a6451854ab7e4b4c146bf78a/preprints/Equality-and-rigidity-in-the-spacetime-Penrose-inequality-September-27-2026/README.md}{Manuscript page and citation}; \href{https://github.com/openai/math/blob/main/preprints/Equality-and-rigidity-in-the-spacetime-Penrose-inequality-September-27-2026/paper.pdf}{Latest version}; Statement source: \href{https://github.com/openai/math/blob/adc7f1241b42e322a6451854ab7e4b4c146bf78a/preprints/Equality-and-rigidity-in-the-spacetime-Penrose-inequality-September-27-2026/build/sections/01-setting.tex\#L154}{Theorem 2.5 (thm:rigidity)}; Statement source: \href{https://github.com/openai/math/blob/adc7f1241b42e322a6451854ab7e4b4c146bf78a/preprints/Equality-and-rigidity-in-the-spacetime-Penrose-inequality-September-27-2026/build/sections/01-setting.tex\#L176}{Theorem 2.6 (thm:converse)}; Statement source: \href{https://github.com/openai/math/blob/adc7f1241b42e322a6451854ab7e4b4c146bf78a/preprints/Equality-and-rigidity-in-the-spacetime-Penrose-inequality-September-27-2026/build/three/01-equality.tex\#L36}{Theorem 8.2 (q3:intro:rigidity)}; Statement source: \href{https://github.com/openai/math/blob/adc7f1241b42e322a6451854ab7e4b4c146bf78a/preprints/Equality-and-rigidity-in-the-spacetime-Penrose-inequality-September-27-2026/build/three/01-equality.tex\#L79}{Theorem 8.3 (q3:intro:connected-rigidity)}.
\item OpenAI, \emph{Spacetime Penrose inequalities: enclosing area, charge, and rigidity}, 2026-10-05; citation key \texttt{OAI:Spacetime-Penrose-inequalities-enclosing-area-charge-and-rigidity-October-5-2026}. \href{https://github.com/openai/math/blob/adc7f1241b42e322a6451854ab7e4b4c146bf78a/preprints/Spacetime-Penrose-inequalities-enclosing-area-charge-and-rigidity-October-5-2026/paper.pdf}{PDF (GitHub)}; \href{https://github.com/openai/math/raw/adc7f1241b42e322a6451854ab7e4b4c146bf78a/preprints/Spacetime-Penrose-inequalities-enclosing-area-charge-and-rigidity-October-5-2026/paper.pdf}{PDF (direct)}; \href{https://github.com/openai/math/blob/adc7f1241b42e322a6451854ab7e4b4c146bf78a/preprints/Spacetime-Penrose-inequalities-enclosing-area-charge-and-rigidity-October-5-2026/README.md}{Manuscript page and citation}; \href{https://github.com/openai/math/blob/main/preprints/Spacetime-Penrose-inequalities-enclosing-area-charge-and-rigidity-October-5-2026/paper.pdf}{Latest version}; Statement source: \href{https://github.com/openai/math/blob/adc7f1241b42e322a6451854ab7e4b4c146bf78a/preprints/Spacetime-Penrose-inequalities-enclosing-area-charge-and-rigidity-October-5-2026/build/sections/neutral/setup.tex\#L92}{Theorem 2.2 (n3:intro:exterior-inequality)}; Statement source: \href{https://github.com/openai/math/blob/adc7f1241b42e322a6451854ab7e4b4c146bf78a/preprints/Spacetime-Penrose-inequalities-enclosing-area-charge-and-rigidity-October-5-2026/build/sections/charged/sections/02-preliminaries.tex\#L63}{ch:intro:main}; Statement source: \href{https://github.com/openai/math/blob/adc7f1241b42e322a6451854ab7e4b4c146bf78a/preprints/Spacetime-Penrose-inequalities-enclosing-area-charge-and-rigidity-October-5-2026/build/sections/neutral-bridge.tex\#L99}{Corollary 9.4 (bridge:bridge:timelike)}; Statement source: \href{https://github.com/openai/math/blob/adc7f1241b42e322a6451854ab7e4b4c146bf78a/preprints/Spacetime-Penrose-inequalities-enclosing-area-charge-and-rigidity-October-5-2026/build/sections/neutral-bridge.tex\#L42}{Theorem 9.2 (bridge:intro:global-penrose)}; Statement source: \href{https://github.com/openai/math/blob/adc7f1241b42e322a6451854ab7e4b4c146bf78a/preprints/Spacetime-Penrose-inequalities-enclosing-area-charge-and-rigidity-October-5-2026/build/sections/rigidity/sections/00-rigidity-setup.tex\#L63}{Theorem 10.2 (n3:intro:rigidity)}.
\item OpenAI, \emph{The Kerr–Newman Penrose Inequality for Axisymmetric Electrovacuum Exteriors}, 2026-10-05; citation key \texttt{OAI:The-Kerr-Newman-Penrose-Inequality-for-Axisymmetric-Electrovacuum-Exteriors-October-5-2026}. \href{https://github.com/openai/math/blob/adc7f1241b42e322a6451854ab7e4b4c146bf78a/preprints/The-Kerr-Newman-Penrose-Inequality-for-Axisymmetric-Electrovacuum-Exteriors-October-5-2026/paper.pdf}{PDF (GitHub)}; \href{https://github.com/openai/math/raw/adc7f1241b42e322a6451854ab7e4b4c146bf78a/preprints/The-Kerr-Newman-Penrose-Inequality-for-Axisymmetric-Electrovacuum-Exteriors-October-5-2026/paper.pdf}{PDF (direct)}; \href{https://github.com/openai/math/blob/adc7f1241b42e322a6451854ab7e4b4c146bf78a/preprints/The-Kerr-Newman-Penrose-Inequality-for-Axisymmetric-Electrovacuum-Exteriors-October-5-2026/README.md}{Manuscript page and citation}; \href{https://github.com/openai/math/blob/main/preprints/The-Kerr-Newman-Penrose-Inequality-for-Axisymmetric-Electrovacuum-Exteriors-October-5-2026/paper.pdf}{Latest version}; Statement source: \href{https://github.com/openai/math/blob/adc7f1241b42e322a6451854ab7e4b4c146bf78a/preprints/The-Kerr-Newman-Penrose-Inequality-for-Axisymmetric-Electrovacuum-Exteriors-October-5-2026/build/main.tex\#L293}{Theorem 2.2 (m:main)}.
\item OpenAI, \emph{The spacetime Penrose inequality and enclosing area}, 2026-09-27; citation key \texttt{OAI:The-spacetime-Penrose-inequality-and-enclosing-area-September-27-2026}. \href{https://github.com/openai/math/blob/adc7f1241b42e322a6451854ab7e4b4c146bf78a/preprints/The-spacetime-Penrose-inequality-and-enclosing-area-September-27-2026/paper.pdf}{PDF (GitHub)}; \href{https://github.com/openai/math/raw/adc7f1241b42e322a6451854ab7e4b4c146bf78a/preprints/The-spacetime-Penrose-inequality-and-enclosing-area-September-27-2026/paper.pdf}{PDF (direct)}; \href{https://github.com/openai/math/blob/adc7f1241b42e322a6451854ab7e4b4c146bf78a/preprints/The-spacetime-Penrose-inequality-and-enclosing-area-September-27-2026/README.md}{Manuscript page and citation}; \href{https://github.com/openai/math/blob/main/preprints/The-spacetime-Penrose-inequality-and-enclosing-area-September-27-2026/paper.pdf}{Latest version}; Statement source: \href{https://github.com/openai/math/blob/adc7f1241b42e322a6451854ab7e4b4c146bf78a/preprints/The-spacetime-Penrose-inequality-and-enclosing-area-September-27-2026/build/sections/introduction.tex\#L87}{Theorem 1.3 (thm:main)}; Statement source: \href{https://github.com/openai/math/blob/adc7f1241b42e322a6451854ab7e4b4c146bf78a/preprints/The-spacetime-Penrose-inequality-and-enclosing-area-September-27-2026/build/sections/introduction.tex\#L110}{Corollary 1.4 (cor:automatic-timelikeness)}; Statement source: \href{https://github.com/openai/math/blob/adc7f1241b42e322a6451854ab7e4b4c146bf78a/preprints/The-spacetime-Penrose-inequality-and-enclosing-area-September-27-2026/build/sections/introduction.tex\#L120}{Theorem 1.5 (thm:low)}; Statement source: \href{https://github.com/openai/math/blob/adc7f1241b42e322a6451854ab7e4b4c146bf78a/preprints/The-spacetime-Penrose-inequality-and-enclosing-area-September-27-2026/build/sections/weak3.tex\#L1269}{Theorem 7.14 (bridge:bridge:timelike / thm:weak3)}; Statement source: \href{https://github.com/openai/math/blob/adc7f1241b42e322a6451854ab7e4b4c146bf78a/preprints/The-spacetime-Penrose-inequality-and-enclosing-area-September-27-2026/build/sections/weak3.tex\#L1354}{Theorem 7.16 (bridge:intro:global-penrose)}.
\item OpenAI, \emph{The spacetime Penrose inequality with charge and original-data rigidity}, 2026-10-05; citation key \texttt{OAI:The-spacetime-Penrose-inequality-with-charge-and-original-data-rigidity-October-5-2026}. \href{https://github.com/openai/math/blob/adc7f1241b42e322a6451854ab7e4b4c146bf78a/preprints/The-spacetime-Penrose-inequality-with-charge-and-original-data-rigidity-October-5-2026/paper.pdf}{PDF (GitHub)}; \href{https://github.com/openai/math/raw/adc7f1241b42e322a6451854ab7e4b4c146bf78a/preprints/The-spacetime-Penrose-inequality-with-charge-and-original-data-rigidity-October-5-2026/paper.pdf}{PDF (direct)}; \href{https://github.com/openai/math/blob/adc7f1241b42e322a6451854ab7e4b4c146bf78a/preprints/The-spacetime-Penrose-inequality-with-charge-and-original-data-rigidity-October-5-2026/README.md}{Manuscript page and citation}; \href{https://github.com/openai/math/blob/main/preprints/The-spacetime-Penrose-inequality-with-charge-and-original-data-rigidity-October-5-2026/paper.pdf}{Latest version}; Statement source: \href{https://github.com/openai/math/blob/adc7f1241b42e322a6451854ab7e4b4c146bf78a/preprints/The-spacetime-Penrose-inequality-with-charge-and-original-data-rigidity-October-5-2026/build/sections/statements.tex\#L103}{Theorem 2.3 (thm:numerical)}; Statement source: \href{https://github.com/openai/math/blob/adc7f1241b42e322a6451854ab7e4b4c146bf78a/preprints/The-spacetime-Penrose-inequality-with-charge-and-original-data-rigidity-October-5-2026/build/sections/statements.tex\#L139}{Theorem 2.5 (thm:rigidity)}; Statement source: \href{https://github.com/openai/math/blob/adc7f1241b42e322a6451854ab7e4b4c146bf78a/preprints/The-spacetime-Penrose-inequality-with-charge-and-original-data-rigidity-October-5-2026/build/sections/statements.tex\#L183}{Theorem 2.6 (thm:converse)}.
\item Family 260: \href{https://github.com/openai/math/blob/adc7f1241b42e322a6451854ab7e4b4c146bf78a/lean/docs/260.md}{Lean scope}.
\item \href{https://github.com/openai/math/blob/adc7f1241b42e322a6451854ab7e4b4c146bf78a/CONTENTS.md}{OpenAI catalogue entry}.
\end{itemize}
\end{document}
